\chapter{Linear Systems, Determinants, and Inverses}

A threat scorer adds three features with unknown weights,
$s=w_{1}f_{1}+w_{2}f_{2}+w_{3}f_{3}$. Calibration on three labeled examples gives three
linear equations at once:
\[
  2w_{1}+w_{2}-w_{3}=8,\qquad
  -3w_{1}-w_{2}+2w_{3}=-11,\qquad
  -2w_{1}+w_{2}+2w_{3}=-3 .
\]
The answer is $w=(2,3,-1)$, and nothing about it is surprising.

Then the coefficient matrix is replaced by the result of another measurement,
\[
  S=\begin{pmatrix}1&2&3\\ 4&5&6\\ 7&8&9.001\end{pmatrix}.
\]
The second row minus the first gives $(3,3,3)$, while the third row minus the second gives
$(3,3,3.001)$: the three rows are nearly dependent. The weights $w=(1,1,1)$ produce the
right-hand side $(6,\,15,\,24.001)$, so the system is consistent and the answer is exactly
$w=(1,1,1)$. Raise the third measurement by $0.001$, and the answer moves to
$(2,-1,2)$. A relative change in the data of $3.46\times10^{-5}$ produces a relative change
in the answer of $1.41$.

There are four questions here, not one: does the system have an answer, how many answers does
it have, which number decides that without solving anything, and how far does the answer shift
when the data shift. Chapter 7 already supplied rank and nullity for the first. Three are still
missing: a procedure that is written down and can be checked by someone else, one number that
decides invertibility and at the same time yields a formula for the answer, and a measure of
sensitivity; those three are the content of this chapter.

The determinant is usually introduced as the cofactor formula, and then its properties are
proved one by one. We reverse that: definition first, through three properties, uniqueness
proved, the cofactor formula appearing as a consequence. The reason is not style: the cofactor
formula hides the fact that the determinant is the only normalized alternating function, and
that fact is what makes the proofs in this chapter short.

\begin{untukanak}
Three friends weigh three boxes on three different scales. Each friend writes down three
numbers, and from those numbers we want to know the weight of each box. As long as the three
scales measure different things, the answer is easy: every new number supplies one new piece
of information.

Here is the trouble. The second scale turns out to nearly copy the first one. The numbers are
not exactly the same, only slightly different, perhaps because of dust on the pan. The third
scale also nearly copies the other two. Now the three pieces of information have nearly the
same content. We feel we hold three pieces of information, when only one or two of them are
really new.

This is where the strangeness begins. If one number changes by a thousandth, the box weights
we compute can change by almost half. Not because the arithmetic is wrong, but because the
information really is nearly identical. There is one number that can tell us in advance
whether this will happen, before we compute anything at all.

Have you ever asked the same question of three people, received three answers with similar
content, and still felt unsure? In this chapter we will learn why similar answers actually
weaken the information, and how to recognize that from the numbers.
\end{untukanak}

\section{Gaussian elimination and three fates of a system}
\tingkat{basic}

A linear system is written as a single unit, not as a list of separate equations. The way to
do it is to attach the right-hand side column to the coefficient matrix. The solution is then
carried out on that arrangement, and every step replaces one system by another system with the
same solution set.

\kamus{elimination}{eliminasi}{Replacing a question by a simpler question without changing the
answer, like simplifying a fraction without changing its value.}

\begin{definisi}[Augmented matrix and elementary row operations]
For the system $A\mathbf{x}=\mathbf{b}$ with $A$ of size $m\times n$, the \emph{augmented
matrix} is the matrix $\bigl(A\mid\mathbf{b}\bigr)$ of size $m\times(n+1)$. The three
\emph{elementary row operations} are: interchanging two rows; multiplying one row by a
nonzero scalar; adding a multiple of one row to another row. Two matrices are called
\emph{row equivalent} if one can be reached from the other by those operations.\index{elimination}
\end{definisi}

Those operations are chosen because each of them can be undone. Interchanging two rows undoes
itself, multiplying by $c\neq0$ is undone by multiplying by $1/c$, and adding $c$ times row $i$
to row $j$ is undone by subtracting $c$ times row $i$. No information is therefore lost.

\begin{teorema}
If two linear systems have row equivalent augmented matrices, the two systems have the same
solution set.
\end{teorema}

\begin{proof}
It is enough to check one operation, because repeated operations follow. Let $S$ be the
original system and $S'$ the system after one operation. Every equation of $S'$ is a linear
combination of the equations of $S$, so every answer of $S$ also answers $S'$. Conversely,
because the operation can be undone by an operation of the same kind, every equation of $S$ is
a linear combination of the equations of $S'$, so every answer of $S'$ also answers $S$. The two
solution sets therefore contain each other.
\end{proof}

The form we aim at is the staircase form. Every nonzero row starts further to the right than
the nonzero row above it, and all entries below that first position are zero. The column where
the first number of a nonzero row sits is called the pivot column, and that number itself is
called a pivot.

\begin{definisi}[Row echelon form]
A matrix is in \emph{row echelon form} if all zero rows lie at the bottom, and the first
nonzero entry of every nonzero row is located further right than the first nonzero entry of the
row above it. The form is called \emph{reduced} if every pivot equals $1$ and is the only
nonzero entry in its column.\index{echelon}\index{pivot}
\end{definisi}

The reduced form is unique. This fact is used every time we speak of ``the rank computed by
elimination'', and its proof is not as short as it looks.

\begin{teorema}[Uniqueness of the reduced echelon form]
Every matrix is row equivalent to exactly one matrix in reduced row echelon form.
\end{teorema}

\begin{proof}
The proof is skipped here. The statement is proved in full in Hoffman and Kunze,
\emph{Linear Algebra} 2nd edition, section 1.4, by showing that the solution space of
$A\mathbf{x}=\mathbf{0}$ determines the position of the pivots uniquely. What we use in this
chapter is only the consequence: the number of pivots and their positions do not depend on the
order of operations chosen.
\end{proof}

With the reduced form, every question about the number of answers can be settled by counting
pivots. These are the three fates of a system.

\begin{teorema}[Three cases]
Suppose $\bigl(A\mid\mathbf{b}\bigr)$ is row equivalent to the reduced echelon form
$\bigl(R\mid\mathbf{c}\bigr)$, with $A$ of size $m\times n$.
\begin{enumerate}
  \item The system is \emph{inconsistent} exactly when there is a pivot in column $(n+1)$.
  \item If it is consistent, the system has a \emph{unique answer} exactly when every column
        of $A$ contains a pivot.
  \item If it is consistent, the system has \emph{many answers} exactly when some column of $A$
        has no pivot, and the number of free parameters equals the number of columns without a
        pivot.
\end{enumerate}
\end{teorema}

\begin{proof}
A pivot in the last column means there is a row of the form $0=\text{nonzero constant}$, which
cannot possibly be satisfied. If there is no such pivot, every nonzero row gives one equation
expressing a pivot variable in terms of the other variables, so the system is consistent.

If every column contains a pivot, the number of pivots equals $n$, so there is no free variable
and every pivot variable is uniquely determined. If some column has no pivot, the variable in
that column is free: it may take any value, and every choice determines the pivot variables
uniquely. The answer is therefore infinite in number, and the number of free parameters is
exactly the number of columns without a pivot.
\end{proof}

Figure~\ref{fig:tiga-nasib} places those three fates in a single plane. The line $x+y=2$ does
not change; only the second line changes. Left panel: $x-y=0$ cuts it at $(1,1)$, two pivots,
a unique answer. Middle panel: $x+y=0$ is parallel, and the last column of the augmented matrix
contains a pivot. Right panel: $2x+2y=4$ is the same line written with different numbers, one
pivot and one free parameter.

\begin{figure}[htbp]
\centering
\begin{tikzpicture}[scale=0.66,>=Stealth]
  \begin{scope}
    \draw[layergrey!28,step=1] (-2.2,-2.2) grid (2.2,2.2);
    \draw[->,ink] (-2.3,0)--(2.4,0) node[right,font=\scriptsize]{$x$};
    \draw[->,ink] (0,-2.3)--(0,2.4) node[above,font=\scriptsize]{$y$};
    \draw[very thick,layerindigo] (-1.2,3.2)--(3.2,-1.2);
    \draw[very thick,layerteal] (-1.6,-1.6)--(2.2,2.2);
    \fill[accent] (1,1) circle (2.4pt);
    \node[font=\scriptsize,below right] at (1,1) {$(1,1)$};
    \node[font=\scriptsize] at (0,-2.9) {(a) intersecting};
  \end{scope}
  \begin{scope}[xshift=6.4cm]
    \draw[layergrey!28,step=1] (-2.2,-2.2) grid (2.2,2.2);
    \draw[->,ink] (-2.3,0)--(2.4,0) node[right,font=\scriptsize]{$x$};
    \draw[->,ink] (0,-2.3)--(0,2.4) node[above,font=\scriptsize]{$y$};
    \draw[very thick,layerindigo] (-1.2,3.2)--(3.2,-1.2);
    \draw[very thick,accent,dashed] (-2.2,2.2)--(2.2,-2.2);
    \node[font=\scriptsize] at (0,-2.9) {(b) parallel};
  \end{scope}
  \begin{scope}[xshift=12.8cm]
    \draw[layergrey!28,step=1] (-2.2,-2.2) grid (2.2,2.2);
    \draw[->,ink] (-2.3,0)--(2.4,0) node[right,font=\scriptsize]{$x$};
    \draw[->,ink] (0,-2.3)--(0,2.4) node[above,font=\scriptsize]{$y$};
    \draw[very thick,layerindigo] (-1.2,3.2)--(3.2,-1.2);
    \draw[line width=2.6pt,accent,densely dotted] (-1.2,3.2)--(3.2,-1.2);
    \node[font=\scriptsize] at (0,-2.9) {(c) coincident};
  \end{scope}
\end{tikzpicture}
\caption{Three fates of a $2\times2$ system with the first line $x+y=2$ held fixed. Left: the
second line $x-y=0$ cuts it at $(1,1)$, two pivots, a unique answer. Middle: the second line
$x+y=0$ is parallel, the right-hand side column contains a pivot, there is no answer. Right: the
second line $2x+2y=4$ coincides with the first line, one pivot, infinitely many answers.}
\label{fig:tiga-nasib}
\end{figure}

\begin{bacakeras}
Two voices. The first voice names the situation, the second voice names the consequence. Read
them alternately, then swap roles.

One system. Two lines. Three fates.

Intersecting: one point, one answer.

Parallel: they never meet, there is no answer.

Coincident: one line disguised as two, infinitely many answers.

Who decides, before we compute anything? The pivots decide.

Two pivots, two columns: a unique answer.

A pivot in the last column: some row reads zero equals nonzero.

A column without a pivot: one free parameter, and the answers line up in one direction.

Repeat the last line together: count the pivots first, then talk about answers.
\end{bacakeras}

The matrix from Chapter 7 displays two fates at once,
\[
  A=\begin{pmatrix}1&2&1\\ 2&5&3\\ 3&7&4\end{pmatrix},
  \qquad \text{third row}=\text{first row}+\text{second row}.
\]
For $\mathbf{b}=(4,9,13)$ elimination gives
\[
  \left(\begin{array}{ccc|c}1&2&1&4\\ 2&5&3&9\\ 3&7&4&13\end{array}\right)
  \longrightarrow
  \left(\begin{array}{ccc|c}1&0&-1&2\\ 0&1&1&1\\ 0&0&0&0\end{array}\right).
\]
There are two pivots and three columns, so there is one free variable. The answer is
$(2,1,0)+t\,(1,-1,1)$ for $t\in\mathbb{R}$, and the check is immediate:
$A(2,1,0)^{\top}=(4,9,13)^{\top}$ while $A(1,-1,1)^{\top}=\mathbf{0}$, so the direction of that
shift really is the blind direction of the matrix from Chapter 7.

For $\mathbf{b}=(4,9,14)$ the last row becomes $0=1$, and the system is inconsistent. The
coefficient matrix is exactly the same; one number on the right-hand side moves the fate of the
system from ``many answers'' to ``no answer''. The question of consistency therefore cannot be
answered from $A$ alone.

The case that appears most often in practice is $m=n$. There the three fates rest on a single
number, and that number is what we look for in the next section.

\begin{akibat}
Suppose $A$ is of size $n\times n$. The system $A\mathbf{x}=\mathbf{b}$ has a unique answer for
every $\mathbf{b}\in\mathbb{R}^{n}$ exactly when $\operatorname{rank}A=n$.
\end{akibat}

\begin{proof}
If $\operatorname{rank}A=n$, the reduced echelon form of $A$ is $I_{n}$, so elimination on
$\bigl(A\mid\mathbf{b}\bigr)$ always ends with a pivot in every column of $A$ and no pivot in
the last column; the previous theorem gives a unique answer. Conversely, if
$\operatorname{rank}A<n$, some column has no pivot, and for a consistent $\mathbf{b}$ the answer
is not unique. Since the rank is less than $n$, the column space of $A$ does not fill
$\mathbb{R}^{n}$, so some $\mathbf{b}$ lies outside it; for that $\mathbf{b}$ the system is
inconsistent.
\end{proof}

\section{The determinant as an alternating function}
\tingkat{intermediate}

We want one number that answers the question of the preceding corollary without solving any
system. That number must be zero exactly when the columns of $A$ are not linearly independent,
must change sign when two columns are interchanged, and must equal one for the identity matrix.
Those three conditions turn out to determine the number uniquely.

\kamus{determinant}{determinan}{The volume factor: how many times space is stretched or
flattened by the matrix.}

\begin{definisi}[Determinant]
A function $D$ from $n\times n$ matrices to $\mathbb{R}$ is called a \emph{determinant} if
\begin{enumerate}
  \item \emph{multilinear in each row}: for scalars $\alpha,\beta$ and rows
        $\mathbf{u},\mathbf{v}$,
        $D(\dots,\alpha\mathbf{u}+\beta\mathbf{v},\dots)
         =\alpha D(\dots,\mathbf{u},\dots)+\beta D(\dots,\mathbf{v},\dots)$;
  \item \emph{alternating}: $D(\dots,\mathbf{u},\dots,\mathbf{u},\dots)=0$ if two rows are
        equal;
  \item \emph{normalized}: $D(I_{n})=1$.
\end{enumerate}
The value $D(A)$ is written $\det A$.\index{determinant}
\end{definisi}

Multilinearity means the determinant is linear in one row while the other rows are held fixed.
The alternating property is the strongest condition, and its consequences are long.

\begin{lemma}
Interchanging two rows changes the sign of the determinant. If one row is a scalar multiple of
another row, or if there is a zero row, the determinant is zero.
\end{lemma}

\begin{proof}
Suppose two adjacent rows both contain $\mathbf{u}+\mathbf{v}$. Since the two rows are equal,
the value is zero. Multilinearity splits that zero into four terms,
\[
  0=D(\dots,\mathbf{u},\mathbf{u},\dots)+D(\dots,\mathbf{u},\mathbf{v},\dots)
   +D(\dots,\mathbf{v},\mathbf{u},\dots)+D(\dots,\mathbf{v},\mathbf{v},\dots),
\]
and the first and fourth terms vanish. Interchanging two adjacent rows therefore reverses the
sign:
\[
  D(\dots,\mathbf{v},\mathbf{u},\dots)=-D(\dots,\mathbf{u},\mathbf{v},\dots).
\]
An interchange of rows that are not adjacent is achieved by several adjacent interchanges, and
every step reverses the sign.

For a multiple, $D(\dots,c\mathbf{u},\dots,\mathbf{u},\dots)
=c\,D(\dots,\mathbf{u},\dots,\mathbf{u},\dots)=0$ by multilinearity and the alternating
property. A zero row is the case $c=0$ of the same form.
\end{proof}

Now uniqueness. The formula that appears in the middle of the proof is the Leibniz formula, and
that formula will later yield the cofactor expansion as a consequence.

\begin{teorema}[Uniqueness and the Leibniz formula]
There is only one function satisfying the three conditions of the definition, and its value is
\[
  \det A=\sum_{\sigma}\operatorname{sgn}(\sigma)\,
  a_{\sigma(1)1}a_{\sigma(2)2}\cdots a_{\sigma(n)n},
\]
where the sum is taken over all permutations $\sigma$ of $\{1,\dots,n\}$ and
$\operatorname{sgn}(\sigma)=(-1)^{k}$ if $\sigma$ is reached by $k$ interchanges.
\end{teorema}

\begin{proof}
Write row $i$ as $\mathbf{a}_{i}=\sum_{j}a_{ij}\mathbf{e}_{j}$. Multilinearity allows every row
to be pulled out into a sum,
\[
  D(A)=\sum_{j_{1},\dots,j_{n}}a_{1j_{1}}a_{2j_{2}}\cdots a_{nj_{n}}\,
  D(\mathbf{e}_{j_{1}},\dots,\mathbf{e}_{j_{n}}).
\]
Terms with $j_{p}=j_{q}$ for $p\neq q$ vanish because $D$ is alternating, so only
$(j_{1},\dots,j_{n})$ that form a permutation $\sigma$ survive. By the previous lemma,
$D(\mathbf{e}_{\sigma(1)},\dots,\mathbf{e}_{\sigma(n)})
=\operatorname{sgn}(\sigma)D(I_{n})=\operatorname{sgn}(\sigma)$.

For uniqueness, suppose $D$ satisfies the three conditions. The same computation gives
$D(A)=\bigl(\sum_{\sigma}\operatorname{sgn}(\sigma)a_{\sigma(1)1}\cdots
a_{\sigma(n)n}\bigr)D(I_{n})=\det A$. Every function satisfying the three conditions therefore
equals that formula.
\end{proof}

The Leibniz formula is correct but impractical: it sums $n!$ terms, and for $n=12$ there are
$479\,001\,600$ of them. What people use is another form of the same formula.

\begin{akibat}[Cofactor expansion]
For every fixed row $i$,
\[
  \det A=\sum_{j=1}^{n}(-1)^{i+j}a_{ij}M_{ij},
\]
where $M_{ij}$ is the determinant of the $(n-1)\times(n-1)$ submatrix obtained by deleting row
$i$ and column $j$. The number $C_{ij}=(-1)^{i+j}M_{ij}$ is called the \emph{cofactor} of the
entry $a_{ij}$. The same formula holds if the expansion is carried out along a column.
\end{akibat}

\begin{proof}
In the Leibniz formula, group the terms by the value of $j$ at position row $i$. Every term
containing $a_{ij}$ has the form $a_{ij}$ times a sign times the product of $n-1$ other
entries, one from each row other than $i$ and one from each column other than $j$. That product
is exactly the Leibniz term for the minor $M_{ij}$, and its sign is $(-1)^{i+j}$ because moving
row $i$ to the last position requires $n-i$ interchanges and moving column $j$ to the last
position requires $n-j$ interchanges, so the sign is $(-1)^{2n-i-j}=(-1)^{i+j}$. Summing over
$j$ gives the formula as written.

For expansion along a column, we use $\det A=\det A^{\top}$. That equality itself follows from
the Leibniz formula: every term for $A^{\top}$ has the form
$a_{\sigma(1)1}\cdots a_{\sigma(n)n}$ with the index order reversed, and the map
$\sigma\mapsto\sigma^{-1}$ preserves the sign because the number of interchanges is the same.
\end{proof}

\begin{contoh}
For
$A=\begin{pmatrix}2&1&-1\\ -3&-1&2\\ -2&1&2\end{pmatrix}$ expansion along the first row gives the
cofactors
\[
  C_{11}=\begin{vmatrix}-1&2\\ 1&2\end{vmatrix}=-4,\qquad
  C_{12}=-\begin{vmatrix}-3&2\\ -2&2\end{vmatrix}=2,\qquad
  C_{13}=\begin{vmatrix}-3&-1\\ -2&1\end{vmatrix}=-5,
\]
so that $\det A=2(-4)+1(2)+(-1)(-5)=-8+2+5=-1$. The alternating sign is not decoration: without
the sign, the result would be $-8+2-5=-11$, and every conclusion afterwards would be wrong.
\end{contoh}

The determinant is useful for product matrices because it is multiplicative.

\begin{teorema}
For $n\times n$ matrices, $\det(AB)=\det A\,\det B$.
\end{teorema}

\begin{proof}
Fix $A$ and define $f(B)=\det(AB)$. Row $i$ of $AB$ is a linear combination of the rows of $B$
with coefficients from row $i$ of $A$, so $f$ is multilinear in the rows of $B$. If two rows of
$B$ are equal, two rows of $AB$ are also equal, so $f$ is alternating. By uniqueness, every
multilinear alternating function $f$ satisfies $f(B)=\det(B)\,f(I_{n})$. Here $f(I_{n})=\det A$,
so $f(B)=\det B\,\det A$.
\end{proof}

\begin{akibat}
$A$ is invertible if and only if $\det A\neq0$. If it is invertible, $\det(A^{-1})=1/\det A$.
Moreover $\det(cA)=c^{n}\det A$, and the determinant of a triangular matrix is the product of
its diagonal entries.
\end{akibat}

\begin{proof}
If $A$ is invertible, $\det A\,\det(A^{-1})=\det I_{n}=1$, so $\det A\neq0$ and its inverse is
$1/\det A$. If $A$ is not invertible, its columns are linearly dependent by Chapter 7, so one
column is a linear combination of the others. Multilinearity splits $\det A$ into a sum of
determinants with two equal columns, so all of them are zero. The formula $\det(cA)=c^{n}\det A$
is obtained by pulling a factor $c$ out of every row. For a triangular matrix, cofactor
expansion along a row that borders only one diagonal entry, repeated $n$ times, leaves the
product of the diagonal.
\end{proof}

That corollary makes elimination and the determinant meet. Elimination brings $A$ to an upper
triangular matrix $U$ by interchanges and additions of row multiples. Additions do not change
the determinant; interchanges reverse its sign. Hence $\det A=(-1)^{s}\det U$ with $s$ the
number of interchanges, and $\det U$ is the product of the pivots: the determinant comes free
after elimination, and that is how libraries compute it.

The other side of the determinant, just as important, is area. The columns of
$B=\begin{pmatrix}3&1\\ 1&2\end{pmatrix}$ are two arrows, $(3,1)$ and $(1,2)$. The unit square
whose sides are carried by those two arrows turns into a parallelogram, and its area is five,
exactly $\det B=3\cdot2-1\cdot1=5$. Figure~\ref{fig:luas-determinan} shows that mapping. If the
two columns are collinear, the parallelogram flattens into a single segment, its area is zero,
and the determinant is zero as well.

\begin{figure}[htbp]
\centering
\begin{tikzpicture}[scale=0.95,>=Stealth]
  \draw[layergrey!28,step=1] (-0.4,-0.4) grid (4.4,3.4);
  \draw[->,ink] (-0.5,0)--(4.6,0) node[right,font=\scriptsize]{$x$};
  \draw[->,ink] (0,-0.5)--(0,3.5) node[above,font=\scriptsize]{$y$};
  \fill[layerteal!16] (0,0)--(3,1)--(4,3)--(1,2)--cycle;
  \draw[thick,layerteal] (0,0)--(3,1)--(4,3)--(1,2)--cycle;
  \fill[layerorange!22] (0,0)--(1,0)--(1,1)--(0,1)--cycle;
  \draw[thick,layerorange] (0,0)--(1,0)--(1,1)--(0,1)--cycle;
  \draw[very thick,accent,->] (0,0)--(3,1)
    node[below right,font=\scriptsize]{$(3,1)$};
  \draw[very thick,accent,->] (0,0)--(1,2)
    node[above left,font=\scriptsize]{$(1,2)$};
  \foreach \t in {1,2,3,4} {\node[font=\tiny,below] at (\t,-0.02) {\t};}
  \foreach \t in {1,2,3} {\node[font=\tiny,left] at (-0.02,\t) {\t};}
  \node[font=\tiny,fill=white,inner sep=0.5pt] at (0.5,0.5) {1};
  \node[font=\scriptsize,fill=white,inner sep=1.5pt] at (2,1.5) {area $5$};
\end{tikzpicture}
\caption{The determinant as area. The unit square (orange, area $1$) is mapped by
$B=\begin{pmatrix}3&1\\ 1&2\end{pmatrix}$ to the parallelogram (teal) spanned by
$(3,1)$ and $(1,2)$. Its area is $5$, the same as $\det B=3\cdot2-1\cdot1=5$: the determinant
multiplies area by five, and its sign records whether the direction of rotation is reversed.}
\label{fig:luas-determinan}
\end{figure}

\begin{tangan}
Materials: graph paper, scissors, ruler, glue.
\begin{enumerate}
  \item Draw a rectangle three squares by one square. Its area is three squares.
  \item Cut out the rectangle. From the left end, cut off a right triangle with base two squares
        and height one square, then glue that piece onto the right end.
  \item The new shape is a slanted parallelogram. Its area is still three squares, because no
        paper was thrown away and none was added.
  \item Repeat with a slant of one square, then two, then three. Each time compute base times
        height, not the length of the slanted side.
\end{enumerate}
What the reader should see: the shape changes, the area does not. A slant only shifts, it does
not stretch, so the determinant does not grow.
\end{tangan}

\begin{untukanak}
Take a rectangular sheet of paper, then slant it by pushing its top edge sideways without
shortening or lengthening anything. Its height stays the same, its width stays the same, its
area stays the same, even though the shape is no longer a rectangle. Only the direction of
rotation changes: trace the bottom edge and then the right edge starting from the lower left
corner, and the slanted sheet still turns the same way; swap the two edges, and the direction
reverses like a left shoe on a right foot.

In this chapter we will learn why a single number can report both area and orientation, and why
that number becomes zero exactly when the sheet is folded into a line. In your opinion, if the
sheet is folded in half, is there still a way to compute its area with the same number?
\end{untukanak}

\section{Cramer's rule and the adjugate}
\tingkat{intermediate}

The determinant answers whether a unique answer exists. The next question: can that answer be
written as a formula, rather than as the result of steps? For small systems it can, and the
formula is tidy.

\kamus{adjugate}{adjoin}{The cofactor map arranged as a transpose, used to write the inverse as
a single formula.}

\begin{definisi}[Adjugate]
For $A$ of size $n\times n$ with cofactors $C_{ij}$, the \emph{cofactor matrix} is the matrix
with entry $C_{ij}$ at position $(i,j)$, and the \emph{adjugate} $\operatorname{adj}A$ is its
transpose, that is, the matrix with entry $C_{ji}$ at position $(i,j)$.
\index{cofactor}\index{adjugate}
\end{definisi}

Note the reversal of the indices: the cofactor $C_{ij}$ comes from deleting row $i$ and column
$j$, but is placed at position $(j,i)$. This is the most frequent source of error when
computing an adjugate by hand.

\begin{teorema}
For every $n\times n$ matrix, $A\operatorname{adj}A
=(\det A)I_{n}=\operatorname{adj}A\,A$.
\end{teorema}

\begin{proof}
Entry $(i,j)$ of $A\operatorname{adj}A$ is $\sum_{k}a_{ik}C_{jk}$. If $i=j$, that sum is exactly
the cofactor expansion along row $i$, namely $\det A$. If $i\neq j$, that sum is the expansion
along row $j$ of the matrix obtained from $A$ by replacing row $j$ with row $i$. That matrix has
two equal rows, so its determinant is zero. Hence $A\operatorname{adj}A=(\det A)I_{n}$. The
second equality is obtained by the same argument on columns, or from
$\operatorname{adj}(A^{\top})=(\operatorname{adj}A)^{\top}$.
\end{proof}

\begin{akibat}
If $\det A\neq0$, then $A^{-1}=\dfrac{1}{\det A}\operatorname{adj}A$.
\end{akibat}

\begin{proof}
From the previous theorem, $A\bigl(\operatorname{adj}A/\det A\bigr)=I_{n}$ and
$\bigl(\operatorname{adj}A/\det A\bigr)A=I_{n}$. The inverse of a matrix is unique, so that
matrix is $A^{-1}$.
\end{proof}

That formula holds for all $n$, but the price is high: $\operatorname{adj}A$ contains $n^{2}$
cofactors, and every cofactor is an $(n-1)\times(n-1)$ determinant.

With that inverse formula, the answer of a system can be written as a ratio of two determinants.

\begin{teorema}[Cramer's rule]
Suppose $A$ is of size $n\times n$ with $\det A\neq0$, and $\mathbf{b}\in\mathbb{R}^{n}$. The
unique answer of $A\mathbf{x}=\mathbf{b}$ is
\[
  x_{j}=\frac{\det A_{j}}{\det A},\qquad j=1,\dots,n,
\]
where $A_{j}$ is the matrix obtained from $A$ by replacing column $j$ with $\mathbf{b}$.
\index{Cramer}
\end{teorema}

\begin{proof}
Since $\det A\neq0$, the previous corollary gives
$\mathbf{x}=A^{-1}\mathbf{b}=(\operatorname{adj}A)\,\mathbf{b}/\det A$. Entry $j$ of that product
is $\sum_{i}b_{i}C_{ji}/\det A$. The sum $\sum_{i}b_{i}C_{ji}$ is the cofactor expansion along
column $j$ of the matrix whose column $j$ has been replaced by $\mathbf{b}$, namely $A_{j}$.
Hence $x_{j}=\det A_{j}/\det A$.
\end{proof}

\begin{contoh}
Back to the opening system, $A\mathbf{x}=\mathbf{b}$ with $\mathbf{b}=(8,-11,-3)^{\top}$. We
already computed $\det A=-1$, and
\[
  \det A_{1}=\begin{vmatrix}8&1&-1\\ -11&-1&2\\ -3&1&2\end{vmatrix}=-2,\qquad
  \det A_{2}=-3,\qquad \det A_{3}=1 .
\]
So $x_{1}=(-2)/(-1)=2$, $x_{2}=(-3)/(-1)=3$, and $x_{3}=1/(-1)=-1$. Check:
$A(2,3,-1)^{\top}=(4+3+1,\,-6-3-2,\,-4+3-2)^{\top}=(8,-11,-3)^{\top}$.

The adjugate of that matrix, with the cofactors of its first row already computed above, is
\[
  \operatorname{adj}A=\begin{pmatrix}-4&-3&1\\ 2&2&-1\\ -5&-4&1\end{pmatrix},
  \qquad
  A^{-1}=\frac{1}{-1}\operatorname{adj}A
  =\begin{pmatrix}4&3&-1\\ -2&-2&1\\ 5&4&-1\end{pmatrix}.
\]
The first row of $A^{-1}$ times $\mathbf{b}$ gives $4(8)+3(-11)-1(-3)=32-33+3=2$, matching
$x_{1}$.
\end{contoh}

Cramer's rule is beautiful and rarely used for computing. Its value lies elsewhere: the formula
shows that the answer depends smoothly on $\mathbf{b}$ and on the entries of $A$, and that
dependence is what we measure through cost and sensitivity.

\section{Elementary matrices, inverses, and equivalent conditions}
\tingkat{advanced}

Row elimination can be seen as matrix multiplication. That view turns the question ``is $A$
invertible'' into a question about a list of factors, and that list of factors yields an
algorithm for the inverse.

\kamus{elementary matrix}{matriks elementer}{The identity matrix with one of its steps already
carried out, so that it records a single row operation.}

\begin{definisi}[Elementary matrix]
An \emph{elementary matrix} is a matrix obtained from $I_{n}$ by one elementary row operation.
There are three kinds, matching the three operations: interchanging two rows, multiplying one
row by a nonzero scalar, and adding a multiple of one row to another row.
\end{definisi}

\begin{lemma}
If $E$ is an $m\times m$ elementary matrix and $A$ an $m\times n$ matrix, then $EA$ is the
result of performing the same row operation on $A$. Every elementary matrix is invertible, and
its inverse is an elementary matrix of the same kind.
\end{lemma}

\begin{proof}
Row $i$ of $EA$ is a linear combination of the rows of $A$ with coefficients from row $i$ of
$E$. For $E$ of the third kind, that row is $\mathbf{a}_{i}+c\mathbf{a}_{j}$, exactly as in the
row operation; for the other two kinds the check is equally direct. The inverse is obtained by
undoing the operation: interchanging rows undoes itself, and
$\operatorname{diag}(1,\dots,c,\dots,1)^{-1}=\operatorname{diag}(1,\dots,1/c,\dots,1)$.
\end{proof}

Figure~\ref{fig:tiga-gerakan} runs the three operations on the system $x+2y=4$, $3x+4y=10$,
which intersect at $(2,1)$. Interchanging rows only interchanges names; multiplying the first
row by $2$ writes the same line with numbers twice as large; adding the second row to the first
gives $4x+6y=14$, which still passes through $(2,1)$. The intersection point does not move, but
the determinant of the coefficient matrix changes sign, doubles, then stays fixed. The solution
set is preserved; the number that measures area is not.

\begin{figure}[htbp]
\centering
\begin{tikzpicture}[scale=0.58,>=Stealth]
  \begin{scope}
    \draw[layergrey!28,step=1] (-0.4,-0.4) grid (4.4,3.2);
    \draw[->,ink] (-0.5,0)--(4.6,0) node[right,font=\scriptsize]{$x$};
    \draw[->,ink] (0,-0.5)--(0,3.4) node[above,font=\scriptsize]{$y$};
    \draw[very thick,layerindigo] (0,2)--(4,0) node[pos=0.15,above,font=\scriptsize]{$R_{1}$};
    \draw[very thick,layerteal] (0,2.5)--(3.33,0) node[pos=0.85,above left,font=\scriptsize]{$R_{2}$};
    \fill[accent] (2,1) circle (2.6pt);
    \node[font=\scriptsize,below right] at (2,1) {$(2,1)$};
    \node[font=\scriptsize] at (2,-1.1) {(i) interchange rows};
  \end{scope}
  \begin{scope}[xshift=5.6cm]
    \draw[layergrey!28,step=1] (-0.4,-0.4) grid (4.4,3.2);
    \draw[->,ink] (-0.5,0)--(4.6,0) node[right,font=\scriptsize]{$x$};
    \draw[->,ink] (0,-0.5)--(0,3.4) node[above,font=\scriptsize]{$y$};
    \draw[very thick,layerindigo] (0,2)--(4,0) node[pos=0.15,above,font=\scriptsize]{$R_{1}$};
    \draw[very thick,layerindigo,densely dashed] (0,2)--(4,0);
    \draw[very thick,layerteal] (0,2.5)--(3.33,0) node[pos=0.85,above left,font=\scriptsize]{$R_{2}$};
    \fill[accent] (2,1) circle (2.6pt);
    \node[font=\scriptsize,below right] at (2,1) {$(2,1)$};
    \node[font=\scriptsize] at (2,-1.1) {(ii) multiply a row};
  \end{scope}
  \begin{scope}[xshift=11.2cm]
    \draw[layergrey!28,step=1] (-0.4,-0.4) grid (4.4,3.2);
    \draw[->,ink] (-0.5,0)--(4.6,0) node[right,font=\scriptsize]{$x$};
    \draw[->,ink] (0,-0.5)--(0,3.4) node[above,font=\scriptsize]{$y$};
    \draw[very thick,layergrey,dash pattern=on 3pt off 2pt] (0,2)--(4,0);
    \draw[very thick,layerteal] (0,2.5)--(3.33,0)
      node[pos=0.8,below right,font=\scriptsize]{$R_{2}$};
    \draw[very thick,layerplum] (0,2.33)--(3.5,0)
      node[pos=0.2,above,font=\scriptsize]{$R_{1}+R_{2}$};
    \fill[accent] (2,1) circle (2.6pt);
    \node[font=\scriptsize,below right] at (2,1) {$(2,1)$};
    \node[font=\scriptsize] at (2,-1.1) {(iii) add a multiple};
  \end{scope}
\end{tikzpicture}
\caption{The three elementary row operations on the system $x+2y=4$, $3x+4y=10$. Interchanging
rows only swaps the names $R_{1}$ and $R_{2}$; multiplying $R_{1}$ by $2$ writes the same line;
adding $R_{2}$ to $R_{1}$ rotates the first line into $4x+6y=14$. In all three panels, both
lines pass through the same intersection point, $(2,1)$.}
\label{fig:tiga-gerakan}
\end{figure}

\begin{teorema}
An $n\times n$ matrix $A$ is invertible if and only if $A$ is a product of elementary matrices.
\end{teorema}

\begin{proof}
A product of invertible matrices is also invertible, so one direction is clear.

For the other direction, suppose $A$ is invertible and bring $A$ to reduced echelon form by row
operations, which is the same as multiplying from the left by elementary matrices
$E_{1},\dots,E_{k}$. By the theorem of three cases, $A$ invertible means every column of $A$
contains a pivot, so the reduced form of $A$ is $I_{n}$, since $A$ is square with $n$ pivots.
Hence $E_{k}\cdots E_{1}A=I_{n}$, and therefore $A=E_{1}^{-1}\cdots E_{k}^{-1}$, a product of
elementary matrices.
\end{proof}

That proof is also an algorithm: performing the same operations on $I_{n}$ gives $A^{-1}$, so
$\bigl(A\mid I_{n}\bigr)\to\bigl(I_{n}\mid A^{-1}\bigr)$ produces the inverse without a single
determinant, far cheaper than the adjugate.

We collect what has been proved into one list. This is the list people use every day, and its
content is another face of a single statement.

\begin{teorema}[Equivalent conditions]
For an $n\times n$ matrix $A$, the following statements are equivalent.
\begin{enumerate}
  \item $A$ is invertible.
  \item $A\mathbf{x}=\mathbf{0}$ has only the trivial answer $\mathbf{x}=\mathbf{0}$.
  \item The reduced echelon form of $A$ is $I_{n}$.
  \item $A$ is a product of elementary matrices.
  \item $A\mathbf{x}=\mathbf{b}$ is consistent for every $\mathbf{b}\in\mathbb{R}^{n}$.
  \item $A\mathbf{x}=\mathbf{b}$ has exactly one answer for every $\mathbf{b}\in\mathbb{R}^{n}$.
  \item $\det A\neq0$.
\end{enumerate}
\end{teorema}

\begin{proof}
We chain the implications $(1)\Rightarrow(2)\Rightarrow(3)\Rightarrow(4)\Rightarrow(1)$,
then $(1)\Rightarrow(5)\Rightarrow(6)\Rightarrow(1)$ and $(1)\Leftrightarrow(7)$.

$(1)\Rightarrow(2)$: multiply $A\mathbf{x}=\mathbf{0}$ from the left by $A^{-1}$.
$(2)\Rightarrow(3)$: the kernel of $A$ is only $\{\mathbf{0}\}$, which means nullity zero, so by
rank-nullity in Chapter 7, $\operatorname{rank}A=n$; hence there are $n$ pivots and the reduced
form is $I_{n}$. $(3)\Rightarrow(4)$: the elimination steps give $E_{k}\cdots E_{1}A=I_{n}$.
$(4)\Rightarrow(1)$: a product of invertible matrices is also invertible.
$(1)\Rightarrow(5)$: take $\mathbf{x}=A^{-1}\mathbf{b}$. $(5)\Rightarrow(6)$: suppose some
$\mathbf{b}$ has two different answers $\mathbf{x}\neq\mathbf{y}$. Then
$\mathbf{x}-\mathbf{y}\neq\mathbf{0}$ and $A(\mathbf{x}-\mathbf{y})=\mathbf{0}$, so the columns
of $A$ are linearly dependent and $\operatorname{rank}A<n$. The column space of $A$ therefore
has dimension less than $n$, so some $\mathbf{b}'$ is out of reach, contradicting (5).
$(6)\Rightarrow(1)$: if $A$ is not invertible, its columns are dependent, so some
$\mathbf{x}\neq\mathbf{0}$ satisfies $A\mathbf{x}=\mathbf{0}$, and the answer for
$\mathbf{b}=\mathbf{0}$ is not unique. Finally $(1)\Leftrightarrow(7)$ is the corollary of the
previous section.
\end{proof}

That list explains why four questions that look different, about solutions, echelon form,
factorization, and the determinant, are in fact one question; only the measuring tool changes.
But the list is silent about cost and sensitivity, and those two things decide whether the
answer obtained can be trusted.

\begin{tantangan}
The following system looks ordinary: three equations for three numbers, no conspicuous
duplicate row, so the answer seems to be unique.
\[
  x+y+z=3,\qquad x+2y+3z=6,\qquad 2x+3y+4z=9 .
\]
Find all its answers. How many answers are there, and why is counting equations misleading here?

Hint 1: add the first two equations, then compare the result with the third equation.

Hint 2: the third equation adds no information at all, so only two pieces of information are
independent. Once $z$ is chosen, $y$ and then $x$ are determined; write the answer as one family
with one parameter.

If you get stuck, go back to the section on three fates and count the pivots.
\end{tantangan}

\subsection*{Cost and sensitivity}

Elimination on an $n\times n$ matrix uses about $n^{3}/3$ multiplications: for $n=1000$ about
$3.33\times10^{8}$, which is still fast. Cramer's rule with cofactors uses $n\cdot n!$
multiplications, about $5.75\times10^{9}$ for $n=12$ and above $10^{19}$ for $n=20$; that is
why it is used only for small $n$ or for theoretical purposes.

The subtler matter is sensitivity. Suppose $A$ is invertible and $\mathbf{b}\neq\mathbf{0}$, and
suppose $\mathbf{b}$ changes to $\mathbf{b}+\delta\mathbf{b}$ so that the answer changes to
$\mathbf{x}+\delta\mathbf{x}$.

\begin{teorema}
With the notation above, for any subordinate matrix norm,
\[
  \frac{\lVert\delta\mathbf{x}\rVert}{\lVert\mathbf{x}\rVert}
  \le\kappa(A)\,
  \frac{\lVert\delta\mathbf{b}\rVert}{\lVert\mathbf{b}\rVert},
  \qquad \kappa(A)=\lVert A\rVert\,\lVert A^{-1}\rVert .
\]
\end{teorema}

\begin{proof}
From $A\,\delta\mathbf{x}=\delta\mathbf{b}$ we get $\delta\mathbf{x}=A^{-1}\delta\mathbf{b}$, so
$\lVert\delta\mathbf{x}\rVert\le\lVert A^{-1}\rVert\lVert\delta\mathbf{b}\rVert$. From
$\mathbf{b}=A\mathbf{x}$ we get $\lVert\mathbf{b}\rVert\le\lVert A\rVert\lVert\mathbf{x}\rVert$,
that is, $1/\lVert\mathbf{x}\rVert\le\lVert A\rVert/\lVert\mathbf{b}\rVert$. Multiplying the two
inequalities gives the required result.
\end{proof}

The number $\kappa(A)$ is called the \emph{condition number}. It does not depend on scale:
multiplying $A$ by $c$ does not change it, because $\lVert cA\rVert\lVert (cA)^{-1}\rVert
=|c|\lVert A\rVert\cdot|c|^{-1}\lVert A^{-1}\rVert$. The determinant has no such property, and
this is why we do not use the determinant as a measure of sensitivity. A small determinant only
means a small volume; it can be small because the whole matrix was scaled down, while the
sensitivity does not change.

For the opening system, $\kappa(S)\approx1.01\times10^{5}$. With machine precision around
$10^{-16}$, the digits that can be trusted in the answer are about $10^{-11}$ relative: enough
for this example, exhausted for a worse matrix. Chapter 22 discusses thresholds and how to
recognize a matrix that can no longer be trusted.
\index{condition}

\begin{untukanak}
Imagine a long plank used as a base for measuring the height of a hill. The plank is almost
horizontal, only slightly tilted. If the lower end is lifted by a finger, the other end rises
far more, because the plank rests on a very small slope. A finger in one place becomes a
hand-span in another.

Now imagine the plank is made very steep. Lift its lower end by a finger, and its upper end
rises only a tiny bit. The same work produces two very different effects, only because the base
is different.

Our computations work like that plank. There are arrangements of numbers in which a small change
in the input produces a small change in the answer; there are also arrangements in which a
change the size of a speck of dust produces a completely different answer. What decides is not
how large the numbers are, but the shape of the arrangement.

In this chapter we will learn why one number can predict that before we compute, and why that
number does not change when all the numbers are enlarged or reduced by the same ratio.

If you weigh on a scale whose base is tilted, is the result wrong, or merely untrustworthy?
\end{untukanak}

\section{A worked example}
\tingkat{intermediate}

We work one system in two ways, compute one $4\times4$ determinant, then measure the sensitivity
of two matrices.

\begin{resep}
\begin{enumerate}
  \item Write the augmented matrix $\bigl(A\mid\mathbf{b}\bigr)$: coefficient columns first, the
        right-hand side column last.
  \item Work on the leftmost column. Choose one nonzero entry as the pivot; if the top entry is
        zero, interchange rows first.
  \item Zero out all entries below the pivot: for each row $j$ below it,
        $R_{j}\leftarrow R_{j}-\frac{a_{ji}}{a_{ii}}R_{i}$.
  \item Shift one column to the right and repeat; stop when there are no nonzero rows left, then
        read the fate of the system from the pivot positions.
\end{enumerate}
Warning about the order: the pivots are zeroed from top to bottom, column by column, before back
substitution begins. Reversing that order destroys the pivots already obtained and forces the
work to be repeated.
\end{resep}

The system is $A\mathbf{x}=\mathbf{b}$ with
\[
  A=\begin{pmatrix}2&1&-1\\ -3&-1&2\\ -2&1&2\end{pmatrix},
  \qquad \mathbf{b}=\begin{pmatrix}8\\ -11\\ -3\end{pmatrix}.
\]

First way, elimination, with the steps written in full:
\[
  \left(\begin{array}{ccc|c}2&1&-1&8\\ -3&-1&2&-11\\ -2&1&2&-3\end{array}\right)
  \xrightarrow[R_{3}\leftarrow R_{3}+R_{1}]{R_{2}\leftarrow R_{2}+\tfrac{3}{2}R_{1}}
  \left(\begin{array}{ccc|c}2&1&-1&8\\ 0&\tfrac{1}{2}&\tfrac{1}{2}&1\\
  0&2&1&5\end{array}\right),
\]
then $R_{3}\leftarrow R_{3}-4R_{2}$ gives
\[
  \left(\begin{array}{ccc|c}2&1&-1&8\\ 0&\tfrac{1}{2}&\tfrac{1}{2}&1\\
  0&0&-1&1\end{array}\right).
\]
Back substitution: from the third row $-x_{3}=1$, so $x_{3}=-1$. The second row
$\tfrac{1}{2}x_{2}+\tfrac{1}{2}(-1)=1$ gives $x_{2}=3$. The first row $2x_{1}+3+1=8$ gives
$x_{1}=2$. The result is $\mathbf{x}=(2,3,-1)$. The pivots that appear are $2$, $\tfrac{1}{2}$,
and $-1$, so the product of the pivots is $\det A=2\cdot\tfrac{1}{2}\cdot(-1)=-1$, matching the
cofactor expansion above. Here there was no row interchange, so the sign did not change.

The second way, Cramer's rule, we already ran on the same system in the previous section. The
three replacement determinants give $\det A_{1}=-2$, $\det A_{2}=-3$, and $\det A_{3}=1$, so
$\mathbf{x}=(-2,-3,1)/(-1)=(2,3,-1)$. Two roads, one answer; a check like this is cheap and
should always be carried out.

Now the $4\times4$ determinant. Take
\[
  M=\begin{pmatrix}2&0&1&3\\ 1&4&0&2\\ 3&1&2&1\\ 0&2&1&5\end{pmatrix}.
\]
Expansion along the first row uses four minors, and one of them we compute in full:
\[
  M_{11}=\begin{vmatrix}4&0&2\\ 1&2&1\\ 2&1&5\end{vmatrix}
  =4(10-1)-0(5-2)+2(1-4)=36-6=30 .
\]
In a similar way $M_{12}=15$, $M_{13}=-45$, and $M_{14}=-15$. Hence
\[
  \det M=2(30)-0(15)+1(-45)-3(-15)=60-45+45=60 .
\]
The number of terms summed is $4!=24$, and for $5\times5$ it is already $120$: the growth is
factorial. This is why libraries do not use this method.

Now sensitivity. The $4\times4$ Hilbert matrix is defined by $h_{ij}=1/(i+j-1)$. Its determinant
is exactly $1/6\,048\,000\approx1.65\times10^{-7}$, very small, so one easily concludes that
this matrix is dangerous. Its condition number is $15\,513.74$. The opening matrix had
$\det S=-0.003$ and $\kappa(S)\approx1.01\times10^{5}$, so its determinant is eighteen
thousand times larger, while its sensitivity is about six and a half times worse. The determinant
is not a measure of sensitivity; the condition number is.

The jump test for $S$ uses $\mathbf{x}=(1,1,1)$, so that $\mathbf{b}=(6,\,15,\,24.001)$.
Changing the third component to $24.002$ changes the answer to $(2,-1,2)$. The relative change
in the data is $3.46\times10^{-5}$, the relative change in the answer is
$\sqrt{2}\approx1.414$, and the ratio is $40\,916$: below $\kappa(S)$ as the theorem promises,
but only by a factor of two and a half.

\kertas{Which is more sensitive, $M$ or the $4\times4$ Hilbert matrix? Predict by comparing their
determinants, then compute the condition number of both with \texttt{numpy.linalg.cond}. Compare
also with $10^{-3}I_{4}$, whose determinant is far smaller, then explain why the rule ``small
determinant means sensitive'' fails. Finally, multiply $M$ by $10^{-4}$: its condition number
does not change, while its determinant shrinks by a factor of $10^{-16}$.}

\section{Computational lab}
\tingkat{advanced}

\begin{lab}
\begin{verbatim}
import numpy as np, math

def det_kofaktor(A):
    n = A.shape[0]
    if n == 1:
        return A[0, 0]
    return sum((-1) ** j * A[0, j] *
               det_kofaktor(np.delete(np.delete(A, 0, 0), j, 1))
               for j in range(n))

A = np.array([[2., 1., -1.], [-3., -1., 2.], [-2., 1., 2.]])
b = np.array([8., -11., -3.])
x = np.linalg.solve(A, b)
print("det A    :", np.linalg.det(A))
print("det cof  :", det_kofaktor(A))
print("rank A   :", np.linalg.matrix_rank(A))
print("x        :", x)
print("residual :", np.linalg.norm(A @ x - b))
for j in range(3):
    Aj = A.copy(); Aj[:, j] = b
    print("det A%d   : %.1f  -> x%d = %.6f" % (j + 1, np.linalg.det(Aj),
          j + 1, np.linalg.det(Aj) / np.linalg.det(A)))

M = np.array([[2., 0., 1., 3.], [1., 4., 0., 2.],
              [3., 1., 2., 1.], [0., 2., 1., 5.]])
print("det M    :", det_kofaktor(M), np.linalg.det(M))

H = np.array([[1. / (i + j + 1) for j in range(4)] for i in range(4)])
print("det H4   :", np.linalg.det(H))
print("cond H4  : %.2f" % np.linalg.cond(H))

S = np.array([[1., 2., 3.], [4., 5., 6.], [7., 8., 9.001]])
bs = S @ np.array([1., 1., 1.])
print("det S    :", np.linalg.det(S), " cond S: %.1f" % np.linalg.cond(S))
print("x        :", np.linalg.solve(S, bs))
bs2 = bs.copy(); bs2[2] += 0.001
print("x'       :", np.linalg.solve(S, bs2))
print("cost     :", math.factorial(12) * 12, round(1000 ** 3 / 3))
\end{verbatim}
Actual output:
\begin{verbatim}
det A    : -1.0
det cof  : -1.0
rank A   : 3
x        : [ 2.  3. -1.]
residual : 1.9860273225978185e-15
det A1   : -2.0  -> x1 = 2.000000
det A2   : -3.0  -> x2 = 3.000000
det A3   : 1.0  -> x3 = -1.000000
det M    : 60.0 60.000000000000036
det H4   : 1.6534391534390412e-07
cond H4  : 15513.74
det S    : -0.0030000000000036655  cond S: 101072.3
x        : [1. 1. 1.]
x'       : [ 2. -1.  2.]
cost     : 5748019200 333333333
\end{verbatim}
\end{lab}

Four things must be read from that output.

First, $x=(2,3,-1)$ appears with residual $1.99\times10^{-15}$, which is zero to machine
precision. The residual is computed on purpose: a library answer is never exact, and the residual
is the cheapest way to check whether that answer satisfies its system. For a bad matrix, a small
residual does not guarantee a small error; that is the topic of Chapter 22.

Second, the determinant from the cofactor expansion and the one from the library are both $-1$,
and for $M$ both are $60$. The library does not use cofactors; it uses a factorization and
multiplies the pivots. The difference of $3.6\times10^{-14}$ on $\det M$ is rounding, not a
difference of definition.

Third, $\det H_{4}$ is reported as $1.65\times10^{-7}$ while its exact value is
$1/6\,048\,000$, so the machine number is correct to seven digits. For higher orders, $1/\det$
passes $10^{12}$ quickly, and that is where the determinant stops being useful as a marker.

Fourth, the jump of the answer for $S$ shows up intact: from $(1,1,1)$ to $(2,-1,2)$ only
because one number changed by a thousandth. The number reported for $\det S$ is not exactly
$-0.003$ either; the difference $3.7\times10^{-15}$ appears because $9.001$ is not
represented exactly in binary fractions.

Finally, cost: $5.748\times10^{9}$ multiplications for Cramer's rule at $n=12$ against
$3.33\times10^{8}$ for elimination at $n=1000$. A far larger matrix is handled with far less
work; that is the difference between a cube and a factorial.

\begin{untukdosen}
A teaching order that has worked: elimination first, with the three cases shown on one fixed
coefficient matrix and three different right-hand sides. Students who see consistency change
because of one number in $\mathbf{b}$ usually accept at once that the question about $A$ and the
question about $\mathbf{b}$ must be separated. Only after that does the determinant enter, as the
answer to the question about $A$ alone.

Two typical misconceptions. First, ``a small determinant means the system is unstable''. The
matrix $10^{-6}I$ has determinant $10^{-18}$ and condition number one; conversely $H_{4}$ has
determinant $1.65\times10^{-7}$ and condition $1.5\times10^{4}$. This mistake appears every
year, and the way to dismantle it is to ask students to multiply one matrix by a scalar and then
compare the two numbers. Second, ``Cramer's rule is the way to solve a system''. Theoretically
true; in cost, false. Give one $12\times12$ system and ask them to estimate the cost of both
ways before computing.

A useful class activity: hand out three systems with the same coefficient matrix and right-hand
sides $(4,9,13)$, $(4,9,14)$, and $(0,0,0)$. Ask each group to determine the fate of its system
from the echelon form alone, then to conclude what can and cannot be concluded from $A$ by
itself.

A sign of readiness for advanced work: students can list the seven equivalent conditions
without memorizing the order, and can explain why the condition about the determinant belongs in
the same list as the conditions about pivots.
\end{untukdosen}

\section{Summary}
\tingkat{basic}

\begin{ringkasan}
\begin{itemize}
  \item Elementary row operations do not change the solution set, because every operation can be
        undone. Because the reduced form is unique, the pivot positions determine three fates:
        inconsistent if a pivot appears in the last column, unique if every coefficient column
        has a pivot, many if some column has no pivot.
  \item The determinant is the only multilinear, alternating, normalized function. The Leibniz
        formula and the cofactor expansion are consequences of it, not its definition.
  \item $\det(AB)=\det A\det B$, so $A$ is invertible exactly when $\det A\neq0$; and $\det A$
        equals the product of the pivots after elimination, times the sign of the interchanges.
  \item The adjugate gives $A\operatorname{adj}A=(\det A)I$ and
        $A^{-1}=\operatorname{adj}A/\det A$; Cramer's rule writes $x_{j}=\det A_{j}/\det A$.
  \item Seven conditions for invertibility are equivalent: trivial kernel, reduced form $I_{n}$,
        product of elementary matrices, consistent and unique for all $\mathbf{b}$, and
        $\det A\neq0$.
  \item The cost of elimination is proportional to $n^{3}/3$, the cost of cofactors to
        $n\cdot n!$. Sensitivity is measured by $\kappa(A)=\lVert A\rVert\lVert A^{-1}\rVert$,
        not by $\det A$.
\end{itemize}
\end{ringkasan}

\section{Exercises}
\tingkat{intermediate}

\begin{latihan}
\begin{enumerate}
  \item Solve by elimination, then check with Cramer's rule:
        \[
          \begin{pmatrix}1&2&3\\ 2&5&3\\ 1&0&8\end{pmatrix}\mathbf{x}
          =\begin{pmatrix}9\\ 8\\ 26\end{pmatrix}.
        \]
        Write down the pivots that appear, record how many times rows were interchanged, then
        conclude $\det A$ from the product of the pivots.

  \item Compute $\det M$ for
        $M=\begin{pmatrix}0&1&2&3\\ 1&0&3&2\\ 2&3&0&1\\ 3&2&1&0\end{pmatrix}$
        by cofactor expansion along one row, then compute it with
        \texttt{numpy.linalg.det}. The two results are not exactly equal: the first is exact, the
        second $5.3\times10^{-15}$. Explain why. Afterwards find a nonzero vector
        $\mathbf{v}$ with $M\mathbf{v}=\mathbf{0}$ and conclude $\det M=0$ without any
        expansion.

  \item Prove that if two columns of $A$ are equal, then $\det A=0$, without using
        $\det A=\det A^{\top}$.
        (Hint: use the Leibniz formula, group the terms by pairs of permutations that swap on
        those two columns, and show that their signs are opposite.)

  \item Prove that $\det(\operatorname{adj}A)=(\det A)^{n-1}$ for $n\ge2$.
        (Hint: take determinants in $A\operatorname{adj}A=(\det A)I$ and use
        $\det(cI)=c^{n}$.)

  \item Evaluate this claim: \emph{if $\det A$ is very small, then the answer of $A\mathbf{x}
        =\mathbf{b}$ is very sensitive to changes in $\mathbf{b}$}. Decide true, false, or not
        yet enough evidence, and include a counterexample or an additional reason that makes it
        true.

  \item For $A=\begin{pmatrix}1&1\\ 1&1+\varepsilon\end{pmatrix}$ with
        $\varepsilon>0$, compute $\kappa(A)$ in the $1$-norm as a function of
        $\varepsilon$. Determine the value of $\varepsilon$ that makes $\kappa(A)>10^{6}$, then
        check it with \texttt{numpy.linalg.cond}.

  \item Prove: if $A$ is invertible and $B$ is obtained from $A$ by adding a multiple of one row
        to another row, then $\det B=\det A$ and $B^{-1}$ can be expressed through $A^{-1}$ and
        one elementary matrix.

  \item Suppose $D=\operatorname{diag}(d_{1},\dots,d_{n})$ and $A$ is of size $n\times n$.
        Show that $\det(A+AD)=\det(A)\prod_{i=1}^{n}(1+d_{i})$, then determine the condition on
        $d_{1},\dots,d_{n}$ for $A+AD$ to be invertible. Check your answer with a $3\times3$
        matrix in Python.

  \item Choose a $3\times3$ matrix with small integer entries and $\det A=1$, compute $A^{-1}$
        through the adjugate and through \texttt{numpy.linalg.inv}, then compare the two. Repeat
        for a matrix with $\det A=7$ and explain why the entries of $A^{-1}$ become fractions.
\end{enumerate}

Hints for the three hardest items. Item 4: take determinants of both sides, remember
$\det(cI)=c^{n}$, and note that $\det A$ may be zero, so the conclusion about
$\operatorname{adj}A$ must be read as a matrix, not as an inverse. Item 6: in the $1$-norm,
$\lVert A\rVert_{1}$ is the largest column sum, and for this matrix the second column is the
largest; compute $\lVert A^{-1}\rVert_{1}$ from the $2\times2$ formula and note that the answer
is proportional to $1/\varepsilon$. Item 8: write $A+AD$ as $A(I+D)$ before pulling out the
determinant, and remember that $I+D$ is diagonal, so its determinant is the product of its
diagonal entries.
\end{latihan}

\begin{untukpeneliti}
What is really proved in this chapter, for the field $\mathbb{R}$ and finite dimension: the
uniqueness of the determinant as a normalized alternating $n$-linear function together with the
permutation formula, in Hoffman and Kunze, \emph{Linear Algebra} 2nd edition, sections 5.2 and
5.3, which hold more generally over a commutative ring with identity (section 5.1);
multiplicativity and the equivalence of invertibility with $\det\neq0$, in Axler,
\emph{Linear Algebra Done Right} 4th edition, section 9C, statements 9.49 and 9.50, with a
basis-free determinant of an operator; and the list of seven equivalent conditions, in Anton
11th edition, Theorem 2.3.8.

What is popularized: ``the determinant is volume''. What is proved is that $|\det A|$ equals the
volume of the parallelepiped spanned by its columns; without the absolute value, the claim is
false for a negative determinant. Also wrong numerically: Cramer's rule as a way to solve
systems above $n=11$, where the cost is factorial.

An open question. The sensitivity bound in the theorem above is normwise and worst-case, often
too pessimistic compared with a componentwise analysis. More fundamentally: the maximum growth
factor of Gaussian elimination with partial pivoting is not known for general $n$. The proved
bound is $2^{n-1}$, while numerical experiments give far less, and that gap is open.

Limits of validity. The argument collapses when the field is not a field, because modules do not
always have a basis and elimination cannot always divide. For very large $n$, the determinant
overflows or underflows long before the rank is lost, so libraries report the logarithm of the
determinant. The sensitivity statements also depend on the norm, and no choice is canonical.
\end{untukpeneliti}

\peta{This chapter follows Anton, \emph{Elementary Linear Algebra} 11th edition, sections 1.1
through 1.6 (linear systems, Gauss-Jordan elimination, inverses, elementary matrices) and 2.1
through 2.3 (cofactors, the determinant through row reduction, Cramer's rule, including Theorem
2.3.8 on equivalent conditions). The axiomatic treatment follows Hoffman and Kunze,
\emph{Linear Algebra} 2nd edition, sections 5.2 and 5.3, with section 1.4 for the uniqueness of
the reduced echelon form. The basis-free view is in Axler, \emph{Linear Algebra Done Right} 4th
edition, section 9C. For elimination as factorization and the view $A=CR$, see Strang,
\emph{Introduction to Linear Algebra} 6th edition, Chapter 2 and the note ``Elimination and
Factorization''. The numerical side, namely condition numbers, rounding error, and cost, is
treated in Chapter 22 of this book.}
